\chapter{Análisis exploratorio: calidad del aire y prevalencia de asma}
\label{ch:analisis}

\section{Objetivo y alcance del análisis}
El objetivo de este capítulo es evaluar, de forma exploratoria, si existe evidencia empírica de una asociación
entre los niveles de contaminación atmosférica y la prevalencia de asma en California a escala de condado.
De acuerdo con la intuición epidemiológica y la literatura previa, cabría esperar que una mayor exposición
a contaminantes se relacione con una mayor prevalencia de asma; sin embargo, los resultados obtenidos en
este primer análisis no reflejan un patrón claro y consistente.

El análisis presentado debe interpretarse como una \textit{primera aproximación} (screening) que sirve para:
(i) validar que la integración de fuentes se ha realizado correctamente, (ii) detectar tendencias o anomalías
gruesas, y (iii) orientar iteraciones posteriores más robustas (multianuales, multivariantes y con control de
confusores).

\section{Metodología}
\subsection{Variables y agregación}
Para cada contaminante se resume la exposición mediante una métrica anual a nivel de condado. La variable
de salud es la prevalencia de asma expresada como proporción/porcentaje de población diagnosticada, también
a nivel de condado, empleando el estrato \textit{Total population -- All ages}.

La integración se realiza por condado (unidad geográfica común), lo que implica que el resultado es un
análisis ecológico: se comparan promedios poblacionales y no exposiciones individuales.

\subsection{Medida de asociación y visualización}
Como medida de asociación lineal se utiliza el coeficiente de correlación de Pearson, $r$, calculado entre:
\begin{itemize}
  \item $x$: nivel anual del contaminante en el condado;
  \item $y$: prevalencia de asma en el condado.
\end{itemize}

Para facilitar la interpretación se representan los condados ordenados de mayor a menor contaminación
(barras) y, sobre el mismo eje de condados, se superpone la prevalencia de asma (línea). Esta visualización
permite comprobar si, al incrementar $x$ (de izquierda a derecha), se observa un incremento sistemático de $y$.

\section{Resultados}
\subsection{Material particulado fino (PM$_{2.5}$)}
La Figura~\ref{fig:pm25_allages} muestra la comparación entre PM$_{2.5}$ y prevalencia de asma en el año base.
A simple vista, no se observa un gradiente monotónico: condados con niveles relativamente altos de PM$_{2.5}$
no presentan necesariamente mayores prevalencias, y algunos máximos de prevalencia aparecen en posiciones
intermedias o incluso con contaminación moderada.

\begin{figure}[htbp]
  \centering
  \includegraphics[width=0.95\textwidth]{figures/ranked_counties_2016__PM___2_5____Total_population__All_ages.png}
  \caption{PM$_{2.5}$ por condado (año base) y prevalencia de asma (\textit{All ages}). Los condados se ordenan
  de mayor a menor nivel del contaminante. La ausencia de un patrón creciente claro sugiere que, en este corte
  temporal, la asociación lineal condal es débil o está enmascarada por otros factores.}
  \label{fig:pm25_allages}
\end{figure}

\subsection{Material particulado grueso (PM$_{10}$)}
La Figura~\ref{fig:pm10_allages} repite el análisis para PM$_{10}$. De nuevo, la prevalencia de asma no aumenta
de forma sistemática con la contaminación. Además, se aprecia que ciertos condados pueden actuar como puntos
de alta influencia (p.\,ej., valores extremos de contaminación) que condicionan cualquier medida lineal global.

\begin{figure}[htbp]
  \centering
  \includegraphics[width=0.95\textwidth]{figures/ranked_counties_2016__PM___10____Total_population__All_ages.png}
  \caption{PM$_{10}$ por condado (año base) y prevalencia de asma (\textit{All ages}). Al igual que con PM$_{2.5}$,
  no se observa un gradiente claro de prevalencia al aumentar el contaminante. La presencia de valores extremos
  sugiere sensibilidad a outliers y refuerza la necesidad de análisis robustos y multivariantes.}
  \label{fig:pm10_allages}
\end{figure}

\section{Discusión: por qué los resultados no se corresponden con lo esperado}
Que no aparezca una correlación clara en estas figuras no implica que la contaminación no tenga efecto sobre
la salud respiratoria; más bien sugiere que, con el diseño actual, el efecto puede quedar \textit{diluido} o
\textit{confundido}. Entre las explicaciones plausibles destacan:

\subsection{Limitaciones temporales y desajustes de periodicidad}
\begin{itemize}
  \item \textbf{Corte en un solo año.} La relación contaminación--asma puede requerir acumulación de exposición,
  rezagos (lags) o cambios sostenidos durante varios años. Un único año reduce drásticamente la potencia para
  detectar señales consistentes.
  \item \textbf{Prevalencia de asma como estimación bianual.} Si la prevalencia se reporta por periodos (p.\,ej.,
  $[year\_from,\,year\_to]$), el valor asignado a un año concreto representa una estimación agregada, lo que
  introduce suavizado temporal y puede ocultar variaciones interanuales.
\end{itemize}

\subsection{Incertidumbre de la variable de salud}
La prevalencia de asma suele estimarse mediante encuestas y/o procedimientos censales indirectos, por lo que
incluye error muestral. Si el análisis utiliza únicamente el punto central, se ignora la variabilidad asociada.
En escenarios con diferencias pequeñas entre condados, esa incertidumbre puede dominar la señal.

\subsection{Heterogeneidad espacial y falacia ecológica}
\begin{itemize}
  \item \textbf{Agregación a nivel condado.} Dentro de un condado pueden coexistir zonas urbanas y rurales con
  exposiciones muy distintas. El promedio condal puede ser un mal proxy de la exposición real de la población.
  \item \textbf{Falacia ecológica.} La ausencia de correlación entre promedios agregados no invalida una relación
  causal a nivel individual, especialmente si existen subpoblaciones vulnerables (infancia, comorbilidades).
\end{itemize}

\subsection{Confusión por factores externos a la calidad del aire}
La prevalencia de asma puede depender de múltiples determinantes que no se han incluido en este análisis:
nivel socioeconómico, acceso y uso del sistema sanitario, hábitos (p.\,ej., tabaquismo), condiciones de vivienda
(humedad, moho), ocupación, alérgenos, clima/temperatura, incendios estacionales, etc. Si estos factores varían
espacialmente y correlacionan con los contaminantes, pueden enmascarar (o incluso invertir) la asociación bivariante.

\subsection{Cobertura y calidad de las mediciones de contaminantes}
La cobertura espacial/temporal de estaciones no es homogénea entre contaminantes ni entre regiones. Condados con
menos sensores pueden presentar estimaciones de exposición más ruidosas o sesgadas, lo que reduce la capacidad de
detectar relaciones. Por tanto, estos resultados deben leerse junto con la evaluación de cobertura de sensores.

\subsection{Relaciones no lineales y efectos de umbral}
Es posible que la relación entre exposición y prevalencia no sea lineal (p.\,ej., umbrales regulatorios, saturación,
o interacción con temperatura/humedad). En esos casos, $r$ de Pearson y una inspección visual ordenada pueden no ser
suficientes para revelar la dinámica real.

\section{Líneas de trabajo futuro}
A partir de estos hallazgos preliminares, las siguientes extensiones son prioritarias:
\begin{itemize}
  \item \textbf{Análisis multianual (serie temporal).} Ampliar el periodo de estudio y modelar tendencias y rezagos
  (p.\,ej., medias móviles, correlaciones con desfase, modelos con efectos temporales).
  \item \textbf{Modelos multivariantes.} Incorporar covariables socio-demográficas y contextuales para controlar
  confusores y estimar efectos parciales.
  \item \textbf{Estratificación por edad.} Repetir el análisis en grupos etarios (infancia/adultos) para detectar
  señales específicas de subpoblaciones vulnerables, cuidando el tamaño muestral efectivo.
  \item \textbf{Componentes espaciales.} Evaluar autocorrelación espacial y considerar modelos con estructura espacial
  (p.\,ej., efectos aleatorios por región, suavizado espacial).
  \item \textbf{Robustez.} Usar medidas resistentes a outliers (Spearman, regresión robusta) y análisis de sensibilidad
  excluyendo condados extremos o con baja cobertura de sensores.
\end{itemize}

En conjunto, este capítulo confirma que, con el corte temporal y el nivel de agregación actuales, no se obtiene
evidencia clara de una correlación lineal fuerte entre contaminación y prevalencia de asma. Lejos de invalidar la
hipótesis, esto justifica la necesidad de un enfoque más rico (temporal, multivariante y espacial) para aproximar
la complejidad del fenómeno.
